Kurs:Algebraische Kurven (Osnabrück 2025-2026)/Arbeitsblatt 20/latex
Kurs:Algebraische Kurven (Osnabrück 2025-2026)/Arbeitsblatt 20/latex

\setcounter{section}{20}


  \zwischenueberschrift{Übungsaufgaben}


Die folgenden vier Aufgaben beweise man einerseits direkt und andererseits mit der Normalität von $\Z$.


\inputaufgabegibtloesung

{}

{


Zeige, dass $\sqrt{2}$ eine
\definitionsverweis {irrationale Zahl}{}{}
ist.


}

{} {}


\inputaufgabe

{}

{


Es sei $p$ eine
\definitionsverweis {Primzahl}{}{.}
Zeige unter Verwendung der
eindeutigen Primfaktorzerlegung
von natürlichen Zahlen, dass die
\definitionsverweis {reelle Zahl}{}{}


\mathl{\sqrt{p}}{}
\definitionsverweis {irrational}{}{}
ist.


}

{} {}


\inputaufgabegibtloesung

{}

{


Beweise mithilfe der eindeutigen Primfaktorzerlegung in $\Z$, dass $9^{1/3}$ irrational ist.


}

{} {}


\inputaufgabegibtloesung

{}

{


Es sei


\mavergleichskette

{\vergleichskette

{ n
}

{ = }{ p_1^{\alpha_1 } \cdots p_r^{\alpha_r}
}

{  }{
}

{    }{
}

{   }{
}

}
{}{}{}
die kanonische Primfaktorzerlegung der natürlichen Zahl $n$. Es sei $k$ eine positive natürliche Zahl und sei vorausgesetzt, dass nicht alle Exponenten $\alpha_i$ ein Vielfaches von $k$ sind. Zeige, dass
die reelle Zahl


\mathdisp {n^{\frac{1}{k} }} {  }


irrational ist.


}

{} {}


\inputaufgabegibtloesung

{}

{


Es sei $R$ ein
\definitionsverweis {normaler Integritätsbereich}{}{}
und

\mathbed {f \in R} {}

 {f \neq 0} {}

 {} {} {} {.} Zeige, dass die Nenneraufnahme $R_f$ ebenfalls normal ist.


}

{} {}


\inputaufgabe

{}

{


Es sei $R$ ein
\definitionsverweis {normaler Integritätsbereich}{}{}
und sei


\mavergleichskette

{\vergleichskette

{S
}

{ \subseteq }{ R
}

{  }{
}

{    }{
}

{   }{
}

}
{}{}{}
ein
\definitionsverweis {multiplikatives System}{}{.}
Zeige, dass dann auch die
\definitionsverweis {Nenneraufnahme}{}{}
$R_S$ normal ist.


}

{} {}


\inputaufgabe

{}

{


Es sei $K$ ein
\definitionsverweis {Körper}{}{}
und sei \mathind { R_i \subseteq K  } { i \in I  }{,} eine Familie von
\definitionsverweis {normalen}{}{}
Unterringen. Zeige, dass auch der Durchschnitt $\bigcap_{i \in I} R_i$ normal ist.


}

{} {}


\inputaufgabe

{}

{


Es sei $R$ ein
\definitionsverweis {Integritätsbereich}{}{.}
Zeige, dass $R$ genau dann
\definitionsverweis {normal}{}{}
ist, wenn er mit seiner
\definitionsverweis {Normalisierung}{}{}
übereinstimmt.


}

{} {}


\inputaufgabe

{}

{


Es sei $R$ ein
\definitionsverweis {Integritätsbereich}{}{.} Es sei angenommen, dass die
\definitionsverweis {Normalisierung}{}{} von $R$ gleich dem
\definitionsverweis {Quotientenkörper}{}{} $Q(R)$ ist. Zeige, dass dann $R$ selbst schon ein
\definitionsverweis {Körper}{}{} ist.


}

{} {}


\inputaufgabe

{}

{


Es sei $M$ ein
\definitionsverweis {torsionsfreies}{}{}
Monoid. Zeige, dass dann auch die
\definitionsverweis {Differenzengruppe}{}{}
$\Gamma(M)$ torsionsfrei ist.


}

{} {}


\inputaufgabe

{}

{


Es sei $M$ eine
\definitionsverweis {kommutative Gruppe}{}{.}
Zeige, dass die
\definitionsverweis {Torsionsfreiheit}{}{}
von $M$ äquivalent zu folgender Eigenschaft ist: Aus


\mavergleichskette

{\vergleichskette

{ m
}

{ \in }{ M
}

{  }{
}

{    }{
}

{   }{
}

}
{}{}{}
und


\mavergleichskette

{\vergleichskette

{ rm
}

{ = }{ 0
}

{  }{
}

{    }{
}

{   }{
}

}
{}{}{}
für ein positives


\mavergleichskette

{\vergleichskette

{ r
}

{ \in }{ \N
}

{  }{
}

{    }{
}

{   }{
}

}
{}{}{}
folgt stets


\mavergleichskette

{\vergleichskette

{ m
}

{ = }{ 0
}

{  }{
}

{    }{
}

{   }{
}

}
{}{}{.}
Zeige ferner, dass diese Äquivalenz für ein
\definitionsverweis {Monoid}{}{}
nicht gelten muss.


}

{} {}


\inputaufgabegibtloesung

{}

{


Wir betrachten das Nullstellengebilde


\mavergleichskettedisp

{\vergleichskette

{ C
}

{ =} { V \left(  Y^2-X^4  \right)
}

{ \subseteq} { {\mathbb A}^{2}_{ {\mathbb C} }
}

{ } {
}

{ } {
}

}
{}{}{.}
\aufzaehlungdrei{Ist $C$ irreduzibel?
}{Kann man den Koordinatenring

\mathl{{\mathbb C}[X,Y]/ \left(  Y^2-X^4  \right)}{} als Monoidring erhalten?
}{Kann man den Koordinatenring

\mathl{{\mathbb C}[X,Y]/ \left(  Y^2-X^4  \right)}{} als Monoidring zu einem Untermonoid


\mavergleichskette

{\vergleichskette

{ M
}

{ \subseteq }{ \N \times \Z/( 2 )
}

{  }{
}

{    }{
}

{   }{
}

}
{}{}{}
erhalten?
}


}

{} {}


\inputaufgabegibtloesung

{}

{


Es sei


\mavergleichskette

{\vergleichskette

{ n
}

{ \in }{ \N_+
}

{  }{
}

{    }{
}

{   }{
}

}
{}{}{}
und sei


\mavergleichskette

{\vergleichskette

{ M
}

{ \subseteq }{ \N \times \Z/( n )
}

{  }{
}

{    }{
}

{   }{
}

}
{}{}{}
ein
\definitionsverweis {Untermonoid}{}{.}
Zeige, dass


\mavergleichskette

{\vergleichskette

{ m
}

{ \in }{ M
}

{  }{
}

{    }{
}

{   }{
}

}
{}{}{}
genau dann eine
\definitionsverweis {Einheit}{}{}
ist, wenn $m$ aufgefasst in

\mathl{\N \times \Z/( n )}{} eine Einheit ist.


}

{} {}


\inputaufgabegibtloesung

{}

{


Es sei


\mavergleichskette

{\vergleichskette

{ n
}

{ \in }{ \N_+
}

{  }{
}

{    }{
}

{   }{
}

}
{}{}{.}
Zeige, dass das
${\mathbb C}$-\definitionsverweis {Spektrum}{}{}
des
\definitionsverweis {kommutativen Monoids}{}{}


\mavergleichskette

{\vergleichskette

{M
}

{ = }{ \N \times \Z/( n )
}

{  }{
}

{    }{
}

{   }{
}

}
{}{}{}
aus $n$
\definitionsverweis {irreduziblen Komponenten}{}{}
besteht, die alle isomorph zur affinen Geraden

\mathl{{\mathbb A}^{1}_{ {\mathbb C} }}{} sind.


}

{} {}


\inputaufgabe

{}

{


Es sei $R$ ein
\definitionsverweis {Integritätsbereich}{}{}
mit
\definitionsverweis {Normalisierung}{}{}
$R ^{\operatorname{norm} }$. Zeige, dass durch


\mavergleichskettedisp

{\vergleichskette

{ {\mathfrak f}
}

{ =} { \left\{  g \in R   \mid   gR^{\operatorname{norm} } \subseteq R         \right\}
}

{ } {
}

{ } {
}

{ } {
}

}
{}{}{}
ein
\definitionsverweis {Ideal}{}{}
in $R$ gegeben ist.


}

{} {}


\inputaufgabe

{}

{


Es sei


\mavergleichskette

{\vergleichskette

{ M
}

{ \subseteq }{ \N
}

{  }{
}

{    }{
}

{   }{
}

}
{}{}{}
ein
\definitionsverweis {numerisches Monoid}{}{,}
das von
\definitionsverweis {teilerfremden}{}{}
natürlichen Zahlen erzeugt werde. Zeige, dass für das
\definitionsverweis {Führungsideal}{}{}
des zugehörigen
\definitionsverweis {Monoidrings}{}{}
$K[M]$ die Beziehung


\mavergleichskettedisp

{\vergleichskette

{ {\mathfrak f}
}

{ =} { (M_{\geq f})
}

{ } {
}

{ } {
}

{ } {
}

}
{}{}{}
besteht, wobei $f$ die
\definitionsverweis {Führungszahl}{}{}
des Monoids bezeichnet.


}

{} {}


\inputaufgabegibtloesung

{}

{


Es sei


\mavergleichskette

{\vergleichskette

{M
}

{ \subseteq }{\N
}

{  }{
}

{    }{
}

{   }{
}

}
{}{}{}
ein
\definitionsverweis {numerisches Monoid}{}{.}
Zeige, dass der
\definitionsverweis {Singularitätsgrad}{}{}
von $M$ mit den drei folgenden Zahlen übereinstimmt.
\aufzaehlungdrei{Die maximale Länge einer Kette von Monoiden


\mavergleichskettedisp

{\vergleichskette

{M
}

{ =} { M_0
}

{ \subset} { M_1
}

{ \subset} { M_2
}

{ \subset  \ldots \subset} { M_n
}

}
{

\vergleichskettefortsetzung

{ =} {\N
}

{ } {}

{ } {}

{ } {}

}{}{}
}{Die maximale Länge einer Kette von
$K$-\definitionsverweis {Algebren}{}{}


\mavergleichskettedisp

{\vergleichskette

{K[M]
}

{ =} { R_0
}

{ \subset} { R_1
}

{ \subset} { R_2
}

{ \subset  \ldots \subset} { R_n
}

}
{

\vergleichskettefortsetzung

{ =} { K[T]
}

{ } {}

{ } {}

{ } {}

}{}{.}
}{Die maximale Länge einer Kette
\zusatzklammer {einer
\definitionsverweis {Fahne}{}{}} {} {}
von
$K$-\definitionsverweis {Untervektor\-räumen}{}{}


\mavergleichskettedisp

{\vergleichskette

{K[M]
}

{ =} { V_0
}

{ \subset} { V_1
}

{ \subset} { V_2
}

{ \subset  \ldots \subset} { V_n
}

}
{

\vergleichskettefortsetzung

{ =} { K[T]
}

{ } {}

{ } {}

{ } {}

}{}{.}
}


}

{} {}


\inputaufgabe

{}

{


Betrachte
Beispiel 20.11.
Welchen Wert haben die drei Erzeuger unter den dort angegebenen Monoidhomomorphismen $\varphi_1,\varphi_2$ nach $\Z$, durch die das Monoid beschrieben werden kann. Bestimme den Kokern des Gruppenhomomorphismus
\maabbeledisp {} {\Gamma(M)  } { \Z^2
} { m } { (\varphi_1(m),\varphi_2(m))
} {.}


}

{(Diesen Kokern nennt man die \stichwort {Divisorenklassengruppe} {} des Monoidringes.)} {}


  \zwischenueberschrift{Aufgaben zum Abgeben}


\inputaufgabe

{3}

{


Es sei $R$ ein
\definitionsverweis {normaler Integritätsbereich}{}{}
und


\mavergleichskette

{\vergleichskette

{ R
}

{ \subseteq }{ S
}

{  }{
}

{    }{
}

{   }{
}

}
{}{}{}
eine
\definitionsverweis {ganze  Ringerweiterung}{}{.}
Es sei


\mavergleichskette

{\vergleichskette

{f
}

{ \in }{ R
}

{  }{
}

{    }{
}

{   }{
}

}
{}{}{.}
Zeige, dass für das von $f$ erzeugte
\definitionsverweis {Hauptideal}{}{}
gilt:


\mavergleichskettedisp

{\vergleichskette

{ R \cap (f)S
}

{ =} { (f)R
}

{ } {
}

{ } {
}

{ } {
}

}
{}{}{.}


}

{} {}


\inputaufgabe

{6}

{


Es sei $R$ ein
\definitionsverweis {normaler Integritätsbereich}{}{.}
Zeige, dass dann auch der
\definitionsverweis {Polynomring}{}{}
$R[X]$ normal ist.


}

{} {}


\inputaufgabe

{5}

{


Es sei $R$ ein
\definitionsverweis {normaler Integritätsbereich}{}{}
und


\mavergleichskette

{\vergleichskette

{ a
}

{ \in }{ R
}

{  }{
}

{    }{
}

{   }{
}

}
{}{}{.}
Es sei vorausgesetzt, dass $a$ keine Quadratwurzel in $R$ besitzt. Zeige, dass das Polynom

\mathl{X^2-a}{}
\definitionsverweis {prim}{}{}
in $R[X]$ ist. Tipp: Verwende den
\definitionsverweis {Quotientenkörper}{}{}
$Q(R)$. Warnung: Prim muss hier nicht zu
\definitionsverweis {irreduzibel}{}{}
äquivalent sein.


}

{} {}


\inputaufgabe

{4}

{


Es sei $R$ ein
\definitionsverweis {Integritätsbereich}{}{.}
Zeige, dass die folgenden Eigenschaften äquivalent sind.
\aufzaehlungdrei{ $R$ ist
\definitionsverweis {normal}{}{.}
}{Für jedes
\definitionsverweis {Primideal}{}{}
${\mathfrak p}$ ist die
\definitionsverweis {Lokalisierung}{}{}
$R_{\mathfrak p}$ normal.
}{Für jedes
\definitionsverweis {maximale Ideal}{}{}
${\mathfrak m}$ ist die Lokalisierung $R_{\mathfrak m}$ normal.}


}

{(Man sagt daher, dass normal eine \stichwort {lokale Eigenschaft} {} ist.)} {}


\inputaufgabe

{2}

{


Es sei


\mavergleichskette

{\vergleichskette

{ M
}

{ \subseteq }{ \Gamma(M)
}

{ \cong }{ \Z^n
}

{    }{
}

{   }{
}

}
{}{}{}
ein
\definitionsverweis {Monoid}{}{}
und betrachte die Menge


\mavergleichskettedisp

{\vergleichskette

{ M^*
}

{ =} { \left\{ \varphi:\Gamma(M) \longrightarrow \Z   \mid  \varphi(M) \subseteq \N         \right\}
}

{ } {
}

{ } {
}

{ } {
}

}
{}{}{.}
Zeige, dass $M^*$ ein normales Untermonoid von

\mathl{\operatorname{Hom} \left(   \Z^n , \Z  \right)}{} ist.


}

{(Dieses Monoid nennt man das \stichwort {duale Monoid} {} zu $M$.)} {}
